\documentclass[11pt]{article}

\usepackage[T1]{fontenc}
\usepackage[utf8]{inputenc}
\usepackage[a4paper,margin=27mm]{geometry}
\usepackage{lmodern}
\usepackage{microtype}
\usepackage{amsmath,amssymb,amsthm,mathtools}
\usepackage{booktabs,tabularx,array}
\usepackage{enumitem}
\usepackage{xcolor}
\usepackage{xurl}
\usepackage{listings}
\usepackage{hyperref}
\usepackage[capitalise,noabbrev]{cleveref}

\definecolor{linkblue}{RGB}{20,70,125}
\definecolor{codegray}{RGB}{247,248,250}
\hypersetup{
  colorlinks=true,
  linkcolor=linkblue,
  citecolor=linkblue,
  urlcolor=linkblue,
  pdftitle={A Machine-Checked Decidability Result for the forall-PO-Free Fragment of ALCI-RCC5},
  pdfauthor={Michael Wessel},
  pdfsubject={Description logics, RCC5, decidability, Lean 4},
  pdfkeywords={description logics, ALCI, RCC5, qualitative spatial reasoning, decidability, Lean, ordered-disjoint structures}
}

\lstdefinelanguage{Lean}{
  morekeywords={def,theorem,structure,inductive,where,let,if,then,else,match,with,fun,forall,exists,by,open,namespace,end,noncomputable},
  sensitive=true,
  morecomment=[l]{--},
  morecomment=[s]{/-}{-/},
  morestring=[b]"
}
\lstset{
  language=Lean,
  basicstyle=\ttfamily\small,
  backgroundcolor=\color{codegray},
  frame=single,
  rulecolor=\color{black!15},
  columns=fullflexible,
  keepspaces=true,
  breaklines=true,
  showstringspaces=false,
  aboveskip=.8em,
  belowskip=.8em
}

\newcolumntype{Y}{>{\raggedright\arraybackslash}X}

\theoremstyle{plain}
\newtheorem{theorem}{Theorem}[section]
\newtheorem{lemma}[theorem]{Lemma}
\newtheorem{proposition}[theorem]{Proposition}
\newtheorem{corollary}[theorem]{Corollary}
\theoremstyle{definition}
\newtheorem{definition}[theorem]{Definition}
\newtheorem{example}[theorem]{Example}
\theoremstyle{remark}
\newtheorem{remark}[theorem]{Remark}

\newcommand{\Logic}{\ensuremath{\mathcal{ALCI}_{\mathrm{RCC5}}}}
\newcommand{\Rel}{\ensuremath{\mathsf{Rel}}}
\newcommand{\EQ}{\ensuremath{\mathsf{EQ}}}
\newcommand{\PP}{\ensuremath{\mathsf{PP}}}
\newcommand{\PPI}{\ensuremath{\mathsf{PPI}}}
\newcommand{\PO}{\ensuremath{\mathsf{PO}}}
\newcommand{\DR}{\ensuremath{\mathsf{DR}}}
\newcommand{\comp}{\ensuremath{\mathsf{comp}}}
\newcommand{\cl}{\ensuremath{\mathsf{cl}}}
\newcommand{\Types}{\ensuremath{\mathsf{Types}}}
\newcommand{\Sig}{\ensuremath{\mathsf{Sig}}}
\newcommand{\Prune}{\ensuremath{\mathsf{Prune}}}
\newcommand{\Sat}{\ensuremath{\mathsf{Sat}}}
\newcommand{\SetSat}{\ensuremath{\mathsf{SetSat}}}
\newcommand{\POFree}{\ensuremath{\mathsf{POFree}}}
\newcommand{\FPOFree}{\ensuremath{\mathsf{FPOFree}}}
\newcommand{\code}[1]{\texttt{#1}}

\setlist[itemize]{leftmargin=2em,itemsep=.25em,topsep=.35em}
\setlist[enumerate]{leftmargin=2em,itemsep=.25em,topsep=.35em}
\setlength{\emergencystretch}{3em}

%% arXiv metadata suggestion:
%%   Primary category: cs.LO
%%   Secondary category: cs.AI
%%   MSC 2020: 03B70, 68T27, 68V15

\title{\textbf{A Machine-Checked Decidability Result for the\\
  $\boldsymbol{\forall\PO}$-Free Fragment of $\boldsymbol{\Logic}$}\\[.45em]
  \large Ordered-Disjoint Models, Finite Cone Schemes,\\
  and Fresh-Occurrence Unfoldings}

\author{Michael Wessel\\
  \small LambdaMikel@Home\\
  \small \texttt{miacwess@gmail.com}}
\date{September 2026}

\begin{document}
\maketitle

\begin{abstract}
We study concept satisfiability for $\Logic$, the description logic
$\mathcal{ALCI}$ whose roles are the five jointly exhaustive, pairwise
disjoint RCC5 relations, under strong equality and the RCC5 converse and
composition laws.  The full decidability problem remains open.  We prove a
sharp unconditional result for the fragment in which negation normal form
contains no universal restriction $\forall\PO.D$.  Existential overlap,
universal and existential disjointness, and both directions of proper-part
modalities remain unrestricted; the fragment does not have the finite-model
property.

The decision procedure is a finite greatest-fixed-point elimination on
\emph{cone signatures}.  A signature records a closure type and a set of types
that may occur strictly below it in the proper-part order.  Completeness is
obtained without a fragment assumption: the signatures realized by any RCC5
model form a post-fixed point of the elimination.  For soundness, each
surviving existential demand is unfolded to a fresh occurrence.  Proper-part
births generate a strict order, disjointness births generate a hereditary
conflict relation, and every remaining pair is interpreted as partial overlap.
This is an ordered-disjoint RCC5 frame by construction.  The absence of
$\forall\PO$ is used exactly where the truth lemma would otherwise have to
control all residual overlap edges.

The complete construction, including normalization of raw formulas, the
polarity-sensitive fragment test, abstract-to-set representation, and an
executable \texttt{Decidable} instance, is formalized in Lean~4 with no
\texttt{sorry}.  The formal capstones decide both abstract satisfiability and
satisfiability by arbitrary nonempty sets.  The explicit signature space is
doubly exponential, so the result is presently a decidability theorem rather
than a practical reasoner.
\end{abstract}

\noindent\textbf{Keywords.}
Description logics; RCC5; qualitative spatial reasoning; modal logic;
decidability; ordered-disjoint structures; greatest fixed points; Lean~4.

\tableofcontents

\section{Introduction}
\label{sec:introduction}

The Region Connection Calculus and the closely related Egenhofer relations
provide a small qualitative vocabulary for spatial configurations.  RCC5
classifies two nonempty regions as equal, a proper part in one direction or
the other, partially overlapping, or disjoint.  Since the five relations are
jointly exhaustive and pairwise disjoint, they form a natural collection of
modal roles.  Cohn proposed modal spatial logics of this kind
\cite{Cohn1993}; Wessel subsequently investigated the $\mathcal{ALCI}_{\rm
RCC}$ family and isolated the RCC5 and RCC8 cases as unresolved
\cite{Wessel2003}.  Lutz and Wolter later studied the corresponding modal
logics systematically and again left the general/substructure RCC5 case open
\cite{LutzWolter2006}.

The difficulty is not merely that $\PP$ is transitive.  Every two domain
elements must receive exactly one RCC5 atom, and every ordered triple must
satisfy the composition table.  Thus a witness created for one modal demand
immediately participates in constraints with all other witnesses.  Moreover,
concepts can force infinite proper-part chains, so a finite-model search is
not complete.  Familiar local blocking arguments are hazardous: identifying
two equally labelled occurrences can create an order cycle or move a
disjointness witness into a cone where downward heredity makes the network
inconsistent.

This paper isolates one polarity boundary at which those global problems can
be controlled.  In negation normal form we forbid only
\[
   \forall\PO.D.
\]
All existential modalities remain, including $\exists\PO.D$, and all
universal modalities over $\EQ,\PP,\PPI,$ and $\DR$ remain.  On raw formulas
the equivalent condition is polarity-sensitive: there is no positive
$\forall\PO$ and no negative $\exists\PO$.  This is necessary because
$\neg\exists\PO.D$ normalizes to $\forall\PO.\neg D$.

\subsection{Contributions}

The main contributions are as follows.

\begin{enumerate}
  \item We give an exact semantic account of strong atomic RCC5 frames as
  \emph{ordered-disjoint structures}: a strict partial order together with a
  symmetric, irreflexive, downward-hereditary disjointness relation.  Partial
  overlap is the residual case.  We also give a canonical representation of
  every such structure by nonempty sets.

  \item We define a finite cone-signature space and a monotone, reductive
  elimination operator.  Iterating the operator for the size of the initial
  space reaches its greatest fixed point.  A concept is accepted when a
  surviving signature contains it.

  \item We prove completeness for \emph{all} concepts.  Given a model, each
  point contributes its closure type and the set of closure types realized by
  its strict proper parts.  The model's realized signatures survive their own
  elimination and therefore survive every round.

  \item We prove soundness for the $\forall\PO$-free fragment by a
  fresh-occurrence unfolding.  No two existential demands share a generated
  element.  Order and disjointness then satisfy the RCC5 frame conditions
  structurally, while the signature conditions propagate every permitted
  universal restriction.

  \item We state the result both for NNF concepts and for formulas with
  unrestricted Boolean negation.  We also prove equivalence between the
  abstract composition-table semantics and a direct semantics over arbitrary
  nonempty sets.

  \item We identify the precise remaining obstruction to using the same
  construction for the full logic: a universal $\PO$ restriction observes
  not only explicit $\PO$ witnesses but every residual pair created by the
  ordered-disjoint completion.
\end{enumerate}

All five steps are represented by named Lean declarations in the accompanying
artifact.  The purpose of this paper is to expose the mathematical spine of
that development, not to reproduce a 45,000-line proof script.

\subsection{Scope and non-claims}
\label{sec:scope}

The result concerns satisfiability of one concept.  It does not cover general
TBoxes, ABoxes, nominals, number restrictions, role inclusions, role
composition expressions in the object language, or the full logic with
$\forall\PO$.  Because the fragment is not closed under negation, decidability
of its satisfiability problem does not automatically yield a same-fragment
subsumption procedure via $C\sqcap\neg D$.  Nor do we claim a tight complexity
classification or an efficient implementation.

\section{Syntax and RCC5 semantics}
\label{sec:semantics}

\subsection{The RCC5 relation algebra}

Let
\[
  \Rel=\{\EQ,\PP,\PPI,\PO,\DR\}.
\]
The converse operation fixes $\EQ,\PO,\DR$ and exchanges $\PP$ with $\PPI$.
Write $\mathcal U=\Rel$.  The composition table used throughout is shown in
\cref{tab:composition}; an entry gives the possible relation between $x$ and
$z$ when the row relation holds from $x$ to $y$ and the column relation from
$y$ to $z$.

\begin{table}[htbp]
\centering
\small
\renewcommand{\arraystretch}{1.18}
\begin{tabular}{c*{5}{c}}
\toprule
$\comp$ & \EQ & \PP & \PPI & \PO & \DR\\
\midrule
\EQ  & $\{\EQ\}$ & $\{\PP\}$ & $\{\PPI\}$ & $\{\PO\}$ & $\{\DR\}$\\
\PP  & $\{\PP\}$ & $\{\PP\}$ & $\mathcal U$ & $\{\PP,\PO,\DR\}$ & $\{\DR\}$\\
\PPI & $\{\PPI\}$ & $\{\EQ,\PP,\PPI,\PO\}$ & $\{\PPI\}$ & $\{\PPI,\PO\}$ & $\{\PPI,\PO,\DR\}$\\
\PO  & $\{\PO\}$ & $\{\PP,\PO\}$ & $\{\PPI,\PO,\DR\}$ & $\mathcal U$ & $\{\PPI,\PO,\DR\}$\\
\DR  & $\{\DR\}$ & $\{\PP,\PO,\DR\}$ & $\{\DR\}$ & $\{\PP,\PO,\DR\}$ & $\mathcal U$\\
\bottomrule
\end{tabular}
\caption{The atomic RCC5 composition table.}
\label{tab:composition}
\end{table}

An \emph{abstract RCC5 frame} is a nonempty set $W$ with a total map
$\rho:W^2\to\Rel$ such that, for all $x,y,z\in W$:
\begin{align*}
  \rho(x,x)&=\EQ,\\
  \rho(x,y)=\EQ&\Longrightarrow x=y,\\
  \rho(y,x)&=\rho(x,y)^{\smile},\\
  \rho(x,z)&\in\comp(\rho(x,y),\rho(y,z)).
\end{align*}
Totality of $\rho$ encodes joint exhaustiveness and functionality encodes
pairwise disjointness of the base relations.  Equality is therefore
\emph{strong}: $\EQ$ is the identity relation, not an equivalence later to be
quotiented.

\subsection{Concepts and satisfaction}

We use the standard negation-normal-form grammar~\cite{BaaderEtAl2007}:
\[
 C,D ::= \top\mid\bot\mid A\mid\neg A\mid C\sqcap D\mid C\sqcup D
        \mid\exists r.C\mid\forall r.C,
 \qquad r\in\Rel.
\]
An interpretation adds a valuation $A^{\mathcal I}\subseteq W$ for every
concept name.  Boolean constructors have their usual meaning and
\begin{align*}
 x\models\exists r.C
   &\quad\Longleftrightarrow\quad
     \text{some }y\in W\text{ has }\rho(x,y)=r\text{ and }y\models C,\\
 x\models\forall r.C
   &\quad\Longleftrightarrow\quad
     \text{every }y\in W\text{ with }\rho(x,y)=r\text{ satisfies }C.
\end{align*}
Concept satisfiability asks whether $x\models C$ for some point of some
abstract RCC5 interpretation.

The formalization permits a carrier type with an explicit domain predicate;
all frame laws and modal quantifiers are relativized to that predicate.  This
is definitionally equivalent to passing to the subtype of in-domain elements,
and nonemptiness is part of satisfiability.

\subsection{The exact fragment}

\begin{definition}[$\forall\PO$-free NNF]
The predicate $\POFree(C)$ is defined structurally.  It recurses through all
Boolean and modal subformulas, permits every existential restriction, and
permits $\forall r.D$ exactly when $r\neq\PO$ and $\POFree(D)$.
\end{definition}

Thus ``$\PO$-free'' would be a misleading abbreviation: overlap itself is
not removed.  The fragment freely contains $\exists\PO.D$ and may mix such
demands with arbitrary nesting of $\PP,\PPI,$ and $\DR$ modalities.

For raw formulas with general negation, define a polarity $b\in\{+,-\}$ while
traversing the syntax.  Negation flips $b$; conjunction and disjunction keep
it.  A positive universal $\PO$ modality and a negative existential $\PO$
modality are forbidden.  Equivalently,
\[
  \FPOFree^+(F) \quad\Longrightarrow\quad
  \POFree(\operatorname{nnf}(F)).
\]
This is the implication used by the decision wrapper.  Notice that
$\neg\forall\PO.A$ is admitted, since its NNF is $\exists\PO.\neg A$,
whereas $\neg\exists\PO.A$ is excluded.

\begin{example}[Composition-forced inconsistency]
Let
\[
 C_{\mathrm{force}}=
   \exists\PP.(\exists\DR.A)\sqcap\forall\DR.\neg A.
\]
If $x\PP y$, $y\DR z$, and $z\models A$, then
$\comp(\PP,\DR)=\{\DR\}$ forces $x\DR z$, contradicting the universal at
$x$.  Hence $C_{\mathrm{force}}$ is unsatisfiable even though it belongs to
the fragment.  A correct procedure must therefore reason with more than the
locally requested edges.
\end{example}

\begin{example}[Failure of the finite-model property]
The fragment concept
\[
 C_{\infty}=\exists\PP.\top\sqcap\forall\PP.\exists\PP.\top
\]
is satisfiable on the natural-number chain, with $m\PP n$ iff $m<n$.
At every $\PP$-successor the universal restriction creates another strict
successor.  Transitivity makes these points pairwise distinct, so no finite
model can satisfy $C_{\infty}$.
\end{example}

\section{Ordered-disjoint normal form}
\label{sec:od}

The key simplification is that the composition table has a compact structural
presentation.

\begin{definition}[Ordered-disjoint structure]
An ordered-disjoint structure is a triple $(N,<,\#)$ such that:
\begin{enumerate}
  \item $<$ is an irreflexive, transitive relation;
  \item $\#$ is symmetric and irreflexive;
  \item $x<y$ implies $\neg(x\#y)$; and
  \item disjointness is downward hereditary in both coordinates:
  \[
      x\#y,\quad x'\leq x,\quad y'\leq y
      \quad\Longrightarrow\quad x'\#y',
  \]
  where $\leq$ is the reflexive closure of $<$.
\end{enumerate}
\end{definition}

It induces an atomic network by priority:
\[
 \rho_{<,\#}(x,y)=
 \begin{cases}
   \EQ  & x=y,\\
   \PP  & x<y,\\
   \PPI & y<x,\\
   \DR  & x\#y,\\
   \PO  & \text{otherwise.}
 \end{cases}
\]
The incompatibility conditions make these cases exclusive.  Partial overlap
is not separately constructed; it is the residual relation left after
identity, comparability, and disjointness have been ruled out.

\begin{proposition}[Ordered-disjoint frame theorem]
\label{prop:odframe}
Every ordered-disjoint structure induces an abstract RCC5 frame.  Conversely,
every abstract RCC5 frame yields an ordered-disjoint structure by defining
$x<y$ iff $\rho(x,y)=\PP$ and $x\#y$ iff $\rho(x,y)=\DR$; the induced network
then recovers $\rho$.
\end{proposition}

\begin{proof}[Proof sketch]
For the forward direction, identity and converse are immediate.  The only
composition cells that force a unique nonidentity output are consequences of
the two structural closure conditions: $\PP;\PP$ and $\PPI;\PPI$ use
transitivity, while $\PP;\DR$ and $\DR;\PPI$ use downward heredity.  In all
other cells the definition of the residual relation and the exclusion axioms
place the output among the listed alternatives.

For the converse, $\comp(\PP,\PP)=\{\PP\}$ gives transitivity.  Symmetry of
$\DR$ follows from converse.  Downward heredity is obtained by two singleton
compositions, first $\PP;\DR=\{\DR\}$ and then
$\DR;\PPI=\{\DR\}$.  Strong equality and the fact that $\rho$ is a function
give irreflexivity and the separation of order from disjointness.  A case split
over the five possible values recovers the original network.
\end{proof}

This normal form is more than convenient notation.  It allows the soundness
construction to prove a small family of geometric invariants rather than all
125 ordered-triple cases directly.

\subsection{Canonical representation by sets}
\label{sec:setrepresentation}

Ordered-disjoint structures also give a direct representation theorem.  Let
$\leq$ be the reflexive closure of $<$.  For each $x\in N$, define a set of
pairs
\[
 \eta(x)=\{(u,v)\in N^2 : \neg(u\#v)\ \text{and}\
                 (u\leq x\ \text{or}\ v\leq x)\}.
\]
The pair $(x,x)$ shows that $\eta(x)$ is nonempty.  The defining axioms imply
\begin{align*}
 \eta(x)\subseteq\eta(y) &\quad\Longleftrightarrow\quad x\leq y,\\
 \eta(x)\cap\eta(y)=\varnothing &\quad\Longleftrightarrow\quad x\#y.
\end{align*}
Antisymmetry of $\leq$ makes $\eta$ injective.  Consequently equality, proper
subset, proper superset, disjointness, and the remaining overlap case between
the sets $\eta(x)$ agree exactly with $\EQ,\PP,\PPI,\DR,$ and $\PO$.

\begin{theorem}[Abstract/set equivalence]
\label{thm:setequiv}
For every $\Logic$ concept $C$, with no fragment assumption,
\begin{align*}
 &C\text{ is satisfiable in an abstract RCC5 frame}\\
 &\qquad\Longleftrightarrow
 C\text{ is satisfiable by a pairwise-distinct family of nonempty sets}
\end{align*}
when RCC5 is interpreted by equality, proper inclusion in either direction,
disjointness, and residual overlap.
\end{theorem}

The right-to-left direction reads an ordered-disjoint structure off an actual
set family.  The left-to-right direction reads an ordered-disjoint structure
off the abstract frame and applies $\eta$.  Both directions, including the
satisfaction transport to the in-domain subtype, are formalized.

\begin{remark}[Regular-closed regions]
\label{rem:regularclosed}
The Lean theorem concerns arbitrary sets; it does not formalize topology.
Lutz and Wolter's representation theorem states that every countable general
RCC5 structure has an isomorphic representation by nonempty regular-closed
sets in $\mathbb R^n$, for every fixed $n>0$ \cite[Theorem~23]{LutzWolter2006}.
The unfolding constructed below is countable, because its occurrences are
finite words over $\Rel\times\mathbb N$.  Combining the two results externally
therefore transfers the fragment theorem to \emph{substructure semantics}
over arbitrary families of regular-closed regions in $\mathbb R^n$.  This does
not establish the corresponding claim for a full structure containing every
regular-closed region: universal modalities are sensitive to adding regions.
\end{remark}

\subsection{Relationship to split-forest models}
\label{sec:splitforest}

The ordered-disjoint normal form is conceptually close to the original
split-forest idea, but the two notions live at different levels.  Both expose
the same semantic backbone: $\PP$ is an order, $\DR$ is hereditary down that
order, and $\PO$ fills the pairs left over.  Moreover, the soundness model used
here really does have a split-forest flavour.  Each non-$\EQ$ existential
demand creates a fresh occurrence; $\PP$ and $\PPI$ births generate vertical
order components, while $\DR$ and $\PO$ births begin new components.  The
result is an occurrence tree whose transitive vertical closure is read as an
ordered-disjoint model.

The difference is substantial.  An ordered-disjoint structure may have an
arbitrary partial order and is an exact semantic normal form, whereas a
split-forest model is a formula-relative presentation obtained by duplicating
or splitting occurrences so that the relevant part of that order has a
tree-like cover with bounded interfaces.  The cone-scheme completeness proof
does not extract such a forest from a given model: it records only a point's
type and the spectrum of types below it.  Tree shape enters later, solely in
the soundness unfolding.  Thus the present method validates the semantic core
of the split-forest intuition while avoiding its hardest full-logic burden,
namely a finite interface that controls every cross-branch relation.  When
$\forall\PO$ is restored, splitting or unravelling creates new residual
$\PO$ pairs that a universal overlap restriction can observe; ordered-
disjointness alone does not show that their target types are coherent.  This
is why the normal form is a natural starting point for future split-forest
work, but is not by itself a full-logic decidability proof.

\section{The finite cone scheme}
\label{sec:scheme}

Fix an NNF concept $C_0$.  Let $\cl(C_0)$ be its list of subformulas, including
$C_0$ itself.  The implementation intentionally permits repeated entries; all
cardinality estimates below are upper bounds in terms of the list length
$m=|\cl(C_0)|$.

\subsection{Types, cones, and static admissibility}

A \emph{candidate type} is any sublist $T$ of $\cl(C_0)$.  Hence at most
$2^m$ candidates are enumerated.  These are deliberately partial Hintikka
types, not maximal consistent sets.  A Boolean support test requires:
\begin{itemize}
  \item $\bot\notin T$ and no pair $A,\neg A$ occurs;
  \item $C\sqcap D\in T$ implies $C,D\in T$;
  \item $C\sqcup D\in T$ implies $C\in T$ or $D\in T$; and
  \item $\forall\EQ.D\in T$ or $\exists\EQ.D\in T$ implies $D\in T$.
\end{itemize}
The last clause uses strong equality: the only $\EQ$-neighbour of a point is
the point itself.

A \emph{cone signature} is a pair
\[
 q=(T,S),
\]
where $T$ is the support type at a point and $S$ is a set of candidate types
allowed strictly below it in the $\PP$ order.  Write
\[
  \operatorname{cone}(q)=\{T\}\cup S.
\]
Besides requiring $T$ and all members of $S$ to pass the support test, static
admissibility requires, for every $U\in S$,
\begin{align}
 \{D:\forall\PP.D\in U\} &\subseteq T,
 \label{eq:sigpp}\\
 \{D:\forall\PPI.D\in T\} &\subseteq U.
 \label{eq:sigppi}
\end{align}
These clauses are the type-level shadow of universal propagation upward and
downward through the proper-part order.

The finite set $S_0=\Sig_0(C_0)$ consists of all statically admissible pairs
drawn from
\[
   \Types(C_0)\times\mathcal P(\Types(C_0)).
\]
Consequently,
\begin{equation}
 |S_0|\leq 2^m\,2^{2^m}=2^{m+2^m}.
 \label{eq:sigbound}
\end{equation}

\subsection{Transition compatibility}

For a demand $\exists r.D\in T$, a target signature
$q'=(T',S')$ must contain $D$ in $T'$.  The additional compatibility
condition depends on $r$ and is summarized in \cref{tab:compat}.  Put
$A_r(T)=\{E:\forall r.E\in T\}$.

\begin{table}[htbp]
\centering
\small
\begin{tabularx}{\textwidth}{@{}cY@{}}
\toprule
Demand & Additional condition on $q=(T,S)$ and $q'=(T',S')$\\
\midrule
$\exists\PP.D$ & $\operatorname{cone}(q)\subseteq S'$\\
$\exists\PPI.D$ & $\operatorname{cone}(q')\subseteq S$\\
$\exists\DR.D$ & for all $U\in\operatorname{cone}(q)$ and
 $V\in\operatorname{cone}(q')$, $A_{\DR}(U)\subseteq V$ and
 $A_{\DR}(V)\subseteq U$\\
$\exists\PO.D$ & $A_{\PO}(T)\subseteq T'$ and
 $A_{\PO}(T')\subseteq T$\\
$\exists\EQ.D$ & handled locally; no transition is generated\\
\bottomrule
\end{tabularx}
\caption{The compatibility test.  In every non-$\EQ$ row, $D\in T'$ is also
required.}
\label{tab:compat}
\end{table}

The $\PP$ clause says that everything at or below the source is strictly below
the target.  Together with \cref{eq:sigpp,eq:sigppi}, this propagates vertical
universals along arbitrarily long chains.  The $\DR$ clause reflects downward
heredity in both coordinates.  The $\PO$ clause propagates direct universal
$\PO$ obligations across an explicit overlap witness in both directions;
although vacuous on the fragment, it is needed for the full-language
completeness theorem discussed in \cref{sec:fullunsat}.

\subsection{Elimination and the decision test}

Given a set $X$ of signatures, define
\[
 \Prune(X)=\bigl\{q\in X :
   \text{every non-$\EQ$ existential demand in $q$ has a compatible target
   in $X$}\bigr\}.
\]
This operator is reductive and monotone:
\[
  \Prune(X)\subseteq X,
  \qquad X\subseteq Y\Longrightarrow\Prune(X)\subseteq\Prune(Y).
\]
Starting with $S_0$, set $S_{i+1}=\Prune(S_i)$.  Since $S_0$ is finite and
each strict change deletes at least one signature, $S_{|S_0|}$ is a fixed
point.  The decision test is
\begin{equation}
  \operatorname{Accept}(C_0)
  \quad\Longleftrightarrow\quad
  \text{some }(T,S)\in S_{|S_0|}\text{ has }C_0\in T.
  \label{eq:accept}
\end{equation}

Equivalently, the executable procedure is:
\begin{enumerate}
  \item compute $\cl(C_0)$ and all candidate types;
  \item enumerate and filter the static cone signatures;
  \item apply $\Prune$ exactly $|S_0|$ times; and
  \item answer \textsc{sat} iff a remaining support type contains $C_0$.
\end{enumerate}
All tests use finite lists and Boolean equality.  The proof obligations are
that no genuine model is eliminated and that every accepted fragment concept
has a model.

\section{Completeness: models form post-fixed points}
\label{sec:completeness}

The completeness proof is short because the finite state space is fixed before
a model is considered.  Let $\mathcal I$ be an arbitrary abstract RCC5 model.
For $x\in W$, define its closure type
\[
  \operatorname{tp}(x)=\{C\in\cl(C_0):x\models C\}
\]
and its strict down-spectrum
\[
  \operatorname{down}(x)=
  \{\operatorname{tp}(y):y\PP x\}.
\]
The model signature of $x$ is
\[
  q_x=(\operatorname{tp}(x),\operatorname{down}(x)).
\]

\begin{lemma}[Static admissibility]
For every in-domain point $x$, $q_x\in S_0$.
\end{lemma}

\begin{proof}[Proof sketch]
Truth immediately gives the Boolean support clauses.  If $y\PP x$ and
$y\models\forall\PP.D$, the actual $\PP$ edge gives $x\models D$.
Conversely, if $x\models\forall\PPI.D$, the same edge read in the inverse
direction gives $y\models D$.  These are exactly the two static vertical
clauses.
\end{proof}

Let
\[
  Z_{\mathcal I}=\{q_x:x\in W\}\cap S_0.
\]
This is a finite set even when the model is infinite, because it is a subset
of the already finite $S_0$.

\begin{lemma}[Model transitions are compatible]
\label{lem:modelcompat}
If $x\models\exists r.D$ and $y$ is a witnessing $r$-neighbour satisfying
$D$, then $q_y$ is compatible with $q_x$ for the demand $(r,D)$.
\end{lemma}

\begin{proof}[Proof sketch]
Certainly $D\in\operatorname{tp}(y)$.  If $x\PP y$, then every point at or
strictly below $x$ is strictly below $y$ by transitivity, giving the $\PP$
cone inclusion.  The $\PPI$ case is dual.  If $x\DR y$, every pair chosen from
the two down-cones is disjoint by downward heredity, so $\DR$ universals cross
in both directions.  If $x\PO y$, symmetry of $\PO$ propagates direct
$\forall\PO$ bodies both ways.  The $\EQ$ case is local because $x=y$.
\end{proof}

\begin{lemma}[Survival]
\label{lem:survival}
$Z_{\mathcal I}\subseteq\Prune(Z_{\mathcal I})$.
\end{lemma}

\begin{proof}
Take $q_x\in Z_{\mathcal I}$ and any non-$\EQ$ existential in its support
type.  The semantics supplies a witness $y$.  By
\cref{lem:modelcompat}, $q_y\in Z_{\mathcal I}$ is a compatible target.
Every demand therefore survives.
\end{proof}

Monotonicity now does all remaining work.  If $Z\subseteq S_0$ and
$Z\subseteq\Prune(Z)$, then induction gives $Z\subseteq S_i$ for every $i$.
In particular every model signature survives the named final round.

\begin{theorem}[Cone-scheme completeness]
\label{thm:complete}
For every NNF concept $C_0$, with no $\POFree$ hypothesis, if $C_0$ is
satisfiable then $\operatorname{Accept}(C_0)$ holds.  More strongly, at every
round $i$ some surviving signature contains $C_0$.
\end{theorem}

\begin{proof}
Choose a model point $x$ satisfying $C_0$.  Then $C_0\in\operatorname{tp}(x)$,
and \cref{lem:survival} plus monotonicity keeps $q_x$ in every $S_i$.
\end{proof}

No bounded-width extraction, filtration quotient, or finite-model property is
used here.  The infinite model is consulted only to establish that one finite
set of its realized signatures is a post-fixed point.

\section{Soundness: the fresh-occurrence unfolding}
\label{sec:soundness}

Assume that $X\subseteq S_0$ is a fixed point in the one-sided sense
$X\subseteq\Prune(X)$, and choose $q_0=(T_0,S_0')\in X$ with $C_0\in T_0$.
For each non-$\EQ$ demand, fix the first compatible target in the finite list
$X$.  This deterministic choice turns $X$ into a finite control graph.

\subsection{Occurrences and births}

An occurrence is a finite word over $\Rel\times\mathbb N$.  The empty word is
the root.  If the current signature's $i$th demand is $(r,D)$, prefixing
$(r,i)$ creates the child that serves it.  Only words generated by this rule
belong to the carrier.  Different words are different domain elements; in
particular, repeated visits to the same finite control state do not identify
occurrences.

This freshness is what allows a finite control graph to describe an infinite
model such as the one required by $C_{\infty}$.  A cycle in the control graph
unfolds to infinitely many distinct words rather than to a forbidden
$\PP$-self-loop.

\subsection{The order}

A $\PP$ birth points upward from a word $w$ to $(\PP,i)::w$; a $\PPI$ birth
points upward in the opposite direction, from $(\PPI,i)::w$ to $w$.  Let $<$
be the transitive closure of these base steps.  Define an integer height by
adding $+1$ for a $\PP$ prefix, $-1$ for a $\PPI$ prefix, and $0$ otherwise.
Every base order step strictly raises height.  Hence $<$ is irreflexive; its
definition makes it transitive.

\subsection{Disjointness}

A $\DR$ birth is a seed pair between $w$ and $(\DR,i)::w$.  Define $\#$ as
the symmetric downward closure of all such seeds with respect to $\leq$.
It is symmetric and downward hereditary by definition.

To see why $<$ and $\#$ cannot collide, define the \emph{vertical base} of a
word by stripping all initial $\PP$ and $\PPI$ births.  Order steps preserve
the vertical base.  A $\DR$ birth changes it, and every pair in $\#$ descends
from the two different bases of such a birth.  Consequently disjoint points
have different vertical bases, while comparable points have the same base.
This proves irreflexivity of $\#$ and $x<y\Rightarrow\neg(x\#y)$.

We have therefore constructed an ordered-disjoint structure on the generated
occurrences.  By \cref{prop:odframe}, it induces a full RCC5 frame.  Each birth
has its advertised relation:
\begin{itemize}
  \item $\PP$ and $\PPI$ births are base order edges;
  \item a $\DR$ birth is a disjointness seed; and
  \item a $\PO$ birth changes vertical base without being a $\DR$ seed, so its
  endpoints are unequal, incomparable, and nondisjoint; it is residual $\PO$.
\end{itemize}

\subsection{Labels and modal propagation}

Label an occurrence with the support type of the control signature reached by
following its word from $q_0$.  Fixed-point membership ensures that every
generated step has a chosen target in $X$, and compatibility ensures that the
target label contains the demanded body.

The universal cases follow the geometry of the unfolding:
\begin{itemize}
  \item Along each proper-part step, compatibility pushes the entire source
  cone into the target's down-spectrum.  This invariant composes along the
  transitive closure.  Static admissibility then propagates $\forall\PP$
  upward and $\forall\PPI$ downward.

  \item Every disjoint pair descends from a $\DR$ birth.  Each endpoint label
  lies in the cone of the corresponding birth endpoint, and the $\DR$
  compatibility clause propagates universal bodies across all such pairs.

  \item $\EQ$ is identity, so its universal and existential bodies are local
  support conditions.

  \item There is no propagation obligation for $\forall\PO$, because no such
  subformula occurs in $\cl(C_0)$ when $\POFree(C_0)$ holds.
\end{itemize}

The last item is the exact point of the fragment restriction.  The generated
frame contains many $\PO$ edges that are not births: every unequal pair that
is neither comparable nor disjoint becomes $\PO$.  Direct transition
compatibility controls explicit $\PO$ witnesses, but it does not control all
of these residual pairs.  A $\forall\PO$ formula would have to inspect every
one of them.

\begin{lemma}[Truth lemma]
\label{lem:truth}
Let $C_0$ be $\forall\PO$-free.  Every concept in the label of a generated
occurrence is true at that occurrence in the unfolded interpretation.
\end{lemma}

\begin{proof}[Proof sketch]
Induct on concepts.  The support test discharges $\bot$, literal clashes,
conjunction, disjunction, and the local $\EQ$ cases.  A non-$\EQ$ existential
uses its own fresh child and the exact birth relation.  The $\PP,\PPI,$ and
$\DR$ universal cases use the propagation facts above.  In the $\PO$
universal case, membership of $\forall\PO.D$ in a label would imply its
membership in $\cl(C_0)$, contradicting $\POFree(C_0)$; the case is therefore
vacuous.
\end{proof}

\begin{theorem}[Cone-scheme soundness]
\label{thm:sound}
If $C_0$ is $\forall\PO$-free and a fixed point $X\subseteq S_0$ contains a
signature whose support type contains $C_0$, then $C_0$ is satisfiable in an
abstract RCC5 frame.  The witnessing frame is countable and is induced by an
ordered-disjoint structure.
\end{theorem}

\begin{proof}
Root the fresh-occurrence unfolding at the chosen signature and apply
\cref{lem:truth} to the empty word.
\end{proof}

\section{Decision theorems}
\label{sec:main}

Combining \cref{thm:complete,thm:sound} at the named round gives the central
characterization.

\begin{theorem}[NNF concept satisfiability]
\label{thm:main}
For every NNF concept $C_0$ with $\POFree(C_0)$,
\[
 C_0\text{ is satisfiable}
 \quad\Longleftrightarrow\quad
 \exists(T,S)\in S_{|S_0|}\; C_0\in T.
\]
Consequently, satisfiability of the $\forall\PO$-free fragment of $\Logic$ is
decidable.
\end{theorem}

Because the right-hand side is a finite Boolean test, the formal result is a
Lean \code{Decidable} value, not merely a classical excluded-middle theorem.

\begin{theorem}[Raw formulas]
\label{thm:raw}
Let $F$ be a formula with unrestricted Boolean negation.  If $F$ has no
positive $\forall\PO$ occurrence and no negative $\exists\PO$ occurrence,
then satisfiability of $F$ is decidable.
\end{theorem}

\begin{proof}[Proof sketch]
The polarity-indexed normalization proves simultaneously that
$\operatorname{nnf}^{+}(F)$ is equivalent to $F$ and
$\operatorname{nnf}^{-}(F)$ is equivalent to $\neg F$.  The structural
polarity condition implies $\POFree(\operatorname{nnf}^{+}(F))$.  Apply
\cref{thm:main} and transport the resulting decision across semantic
equivalence.
\end{proof}

Together with \cref{thm:setequiv}, the same algorithm decides satisfiability
over pairwise-distinct families of arbitrary nonempty sets.

\begin{corollary}[Concrete set satisfiability]
For every $\forall\PO$-free concept $C_0$, satisfiability under the direct set
interpretation of RCC5 is decidable.
\end{corollary}

\subsection{A one-sided theorem for the full language}
\label{sec:fullunsat}

Since completeness never used $\POFree$, it yields a useful but limited
corollary.

\begin{corollary}[Sound full-language refutation]
\label{cor:fullunsat}
For any NNF concept $C_0$, if no signature surviving any chosen elimination
round contains $C_0$, then $C_0$ is unsatisfiable.  In particular this test is
sound at the named round $|S_0|$.
\end{corollary}

The converse is not claimed.  A surviving signature outside the fragment need
not unfold to a model, because residual $\PO$ edges may violate a universal
$\PO$ restriction.  Thus \cref{cor:fullunsat} is a finite one-sided analyzer,
not a decision procedure for full $\Logic$.

\section{Formalization and verification boundary}
\label{sec:formalization}

\subsection{Lean theorem map}

The development uses Lean~4 core and represents all finite control structures
as lists.  \Cref{tab:theoremmap} maps the mathematical results above to the
capstone declarations in \path{formal/POFreeLift.lean} \cite{Artifact2026}.

\begin{table}[htbp]
\centering
\small
\begin{tabularx}{\textwidth}{@{}Yl@{}}
\toprule
Mathematical role & Lean declaration\\
\midrule
Ordered-disjoint structure induces an RCC5 frame & \code{odNet\_frame}\\
Every abstract model reads off as ordered-disjoint & \code{odOfModel}, \code{odNet\_odOfModel}\\
Static signatures and transition test & \code{sigStatic}, \code{compatB}\\
Reductive, monotone elimination and fixed point & \code{pruneSig\_red}, \code{pruneSig\_mono}, \code{coneScheme\_gfp}\\
Model signatures form a post-fixed point & \code{modelSigs\_survives}, \code{modelSigs\_in\_gfp}\\
Full-language completeness & \code{coneScheme\_complete}\\
Generated ordered-disjoint carrier & \code{gOD}, \code{unfInterp\_rcc5}\\
Universal propagation / existential service & \code{allPP\_gLt}, \code{allPPI\_gLt}, \code{allDR\_gDisj}, \code{gchild\_serves}\\
Fragment truth lemma and soundness & \code{unf\_truth}, \code{coneScheme\_sound}\\
Named-round characterization and NNF decision & \code{coneScheme\_correct\_at}, \code{decidableSat\_cone}\\
Raw normalization and decision & \code{nnfP\_correct}, \code{pofree\_nnfP}, \code{decidableFSat}\\
Abstract/set equivalence and set decision & \code{satisfiable\_iff\_set}, \code{decidableSetSat}\\
One-sided full-language refutation & \code{coneScheme\_unsat\_full}, \code{coneScheme\_unsat\_full\_at}\\
\bottomrule
\end{tabularx}
\caption{Result-to-artifact map for the cone-scheme proof.}
\label{tab:theoremmap}
\end{table}

The source contains no \code{sorry}, user-declared axiom, unsafe definition,
\code{native\_decide}, or implementation override in the audited snapshot.
The printed axiom footprint of the main capstones is
\code{propext}, \code{Classical.choice}, and \code{Quot.sound}, the standard
principles used by Lean's classical reasoning and quotient infrastructure.
Classical choice is used in model-side existence proofs and in a stabilization
argument; the finite decision data itself is computed by Boolean list
operations.  In particular, \code{decidableSat\_cone} and
\code{decidableFSat} are executable definitions.  The artifact includes small
kernel-reduced positive and negative examples.

\subsection{Audited snapshot and reproducibility}

This manuscript describes repository commit
\code{795b2707b72e927c56b465ef0413b77aeb9a66a6}.  The relevant files in that
snapshot have the following identifiers:

\begin{center}
\small
\begin{tabularx}{\textwidth}{@{}lY@{}}
\toprule
File & SHA-256\\
\midrule
\path{formal/POFreeLift.lean} &
\code{5f85ff556d0d67205432311f136a903}\allowbreak
\code{35b4a0540576604ec0f395e4aa3db28ae}\\
\path{formal/RCC5NormalForm.lean} &
\code{438e8277891054c82050d2d1c830651}\allowbreak
\code{00881b0630692b72c7c71a5ee80d127ad}\\
\bottomrule
\end{tabularx}
\end{center}

The cumulative \path{POFreeLift.lean} file has 45,375 lines; much of it records
earlier constructions, while the current cone scheme and its capstones occur
near the end.  Recorded independent audits built the development under Lean
4.32.0 and 4.33.1 with zero errors, warnings, and sorries.  The minimal build
command is:
\begin{lstlisting}
lean POFreeLift.lean
\end{lstlisting}
The snapshot does not contain a \path{lean-toolchain} file or a Lake project.
Before archival release, the repository should pin the tested Lean version and
provide a minimal module containing only the theorem's dependency closure.
This is a reproducibility and reviewability improvement, not a logical premise
of the theorem.

\subsection{Adversarial testing}

The formal proof is the correctness evidence; testing is supplementary.  The
repository nevertheless contains independent probes that rederive the RCC5
table from finite set models, compare the finite control layer with exhaustive
small-model search, mutate compatibility conditions, and instantiate bounded
prefixes of the unfolding.  Three staged cold reviews focused respectively on
the trusted semantic base, completeness, and soundness.  The completeness
review found that an early $\PO$ transition clause omitted symmetric
$\forall\PO$ propagation.  The clause was strengthened, and its preservation
of completeness and fragment soundness was then proved in Lean.  The
soundness-focused review traced apparent probe failures to hypotheses that the
probe had abstracted away and led to corresponding probe corrections.  No
review found a counterexample to the fragment capstone.

These experiments should not be conflated with proof, and their finite-model
or bounded-depth nature cannot validate the infinite unfolding.  Their value
is different: they exposed transcription errors, under-strength conditions,
and misleading prose before release.

\section{Complexity and practical status}
\label{sec:complexity}

Let $m=|\cl(C_0)|$.  Candidate types number at most $2^m$, and
\cref{eq:sigbound} gives at most $2^{m+2^m}$ raw signatures before static
filtering and deduplication.  A direct implementation checks, for every
surviving signature and every existential demand, whether some target is
compatible, and performs at most $|S_0|$ rounds.  Polynomially many scans of
an explicitly represented doubly-exponential state space remain within a
coarse doubly-exponential time upper bound.

The Lean development proves finiteness and termination but does not formalize
this asymptotic analysis, nor does it establish a matching lower bound.  More
importantly, the literal list implementation materializes a power set of the
type space.  Audit runs report that only very small closures are feasible and
that a one-modality example can exhaust substantial memory.  The procedure is
therefore not a competitive reasoner.

There is considerable room for symbolic improvement.  The second component of
a signature is used through inclusions and universal-body summaries, suggesting
antichain representations, binary decision diagrams, or SAT/SMT encodings.
Such changes would require new correctness proofs: quotienting two spectra
merely because they agree on currently inspected summaries risks reintroducing
the same joint-coherence problem that freshness was designed to avoid.

\section{Related work}
\label{sec:related}

\paragraph{RCC and qualitative spatial reasoning.}
RCC8 was introduced as a region-based spatial logic by Randell, Cui, and Cohn
\cite{RandellCuiCohn1992}; RCC5 is its coarser five-relation calculus.  Renz
and Nebel studied tractable relation-algebra fragments for RCC reasoning
\cite{RenzNebel1999}.  Those constraint-satisfaction results concern finite
networks of spatial variables.  The present modal satisfiability problem adds
arbitrary nesting, unary concept predicates, and models that may need to be
infinite, so finite RCC network consistency is an ingredient rather than a
decision procedure.

\paragraph{Modal RCC5.}
Cohn proposed modal and non-modal qualitative spatial languages
\cite{Cohn1993}.  Wessel developed the $\Logic$ family and documented both
decidable and undecidable neighbouring systems, leaving the RCC5 case open
\cite{Wessel2003}.  Lutz and Wolter formalized modal logics over RCC8 and RCC5
region structures \cite{LutzWolter2006}.  Their general RCC5 structures have
the same essential conditions used here: JEPD atoms, identity equality,
converse, and the composition table.  Their Theorem~23 supplies the
regular-closed representation used in \cref{rem:regularclosed}.  Their RCC5
undecidability theorem assumes a stronger least-superregion property and
applies to full classes satisfying it; they explicitly leave other RCC5
logics, including the general/substructure case, open.  It therefore does not
subsume the fragment theorem presented here.

The syntactic cut considered here is also different from simply dropping the
$\PO$ modality.  Existential $\PO$ remains available; only universal access to
all $\PO$ neighbours is removed after NNF.  We are not aware of an earlier
published result isolating this one-polarity fragment, but this narrow novelty
claim should be checked independently in peer review.

\paragraph{Description logics with concrete domains.}
Concrete-domain extensions connect abstract individuals to concrete values by
features and predicates.  Tableau and admissibility results for such systems
include Lutz and Mili\v{c}i\'c \cite{LutzMilicic2007} and later work by Baader
and Rydval \cite{BaaderRydval2020}.  They do not directly settle the present
logic: here the region objects themselves are the modal domain and the RCC5
relations connect every pair of them.

\paragraph{Machine-checked metatheory.}
Lean~4 combines dependent type theory with an extensible elaborator and proof
automation \cite{deMouraUllrich2021}.  In this development it serves two
roles: checking the finite fixed-point argument and checking the infinite
model construction.  The latter is the more consequential use, since the
fresh-occurrence carrier, the order/disjointness invariants, and the modal
truth lemma are all quantified over unbounded objects.

\section{\texorpdfstring{Why full $\Logic$ remains open}{Why full ALCI-RCC5 remains open}}
\label{sec:full}

The current method draws a clean line between what is already solved and what
is not.  Completeness works for a full-language input because a real model
supplies compatible witnesses.  Soundness fails to generalize because a
control transition sees only an explicit witness edge, whereas the induced
ordered-disjoint frame assigns $\PO$ to every residual pair.

Consider a source occurrence labelled $\forall\PO.D$.  It is not enough to
ensure $D$ at the target of each explicit $\exists\PO$ birth.  The source may
be residual-$\PO$ related to occurrences created several generations away,
to siblings created for unrelated demands, or to points in other vertical
components.  Whether such a pair is residual depends jointly on the absence of
identity, order in either direction, and hereditary disjointness.  A finite
control state must either predict the types of all those targets or impose a
model shape in which their obligations can be summarized.

This observation suggests four concrete research directions.

\begin{enumerate}
  \item \textbf{Symbolic residual neighbourhoods.}  Enrich a signature with a
  finite summary of which types may be residual-$\PO$ related to its cone, and
  seek a monotone elimination whose post-fixed-point completeness can still be
  read directly from models.

  \item \textbf{Formula-relative split forests.}  Start from the
  ordered-disjoint normal form, split only the occurrences needed to obtain a
  tree-like cover, and carry bounded interfaces that explicitly propagate
  $\forall\PO$ obligations across forgotten branches.  The missing theorem is
  a completeness bound on those interfaces, not the existence of the
  underlying order/disjointness structure.

  \item \textbf{Disjunctive composition information.}  Cone signatures encode
  consequences of singleton-forcing cells particularly well.  Full $\PO$
  control may require finite records of multi-path or disjunctive constraints,
  together with a proof that the resulting operator remains monotone.

  \item \textbf{A negative boundary result.}  If no finite residual
  neighbourhood summary can be made complete, that obstruction may guide an
  undecidability reduction.  Any such reduction must force the coincidence or
  coordination of independently generated witnesses; merely constructing a
  grid-shaped RCC5 model is insufficient.
\end{enumerate}

The ordered-disjoint normal form should precede these attacks.  It identifies
exactly which parts of RCC5 are positive closure conditions and which part is
the complement.  The full problem is, in this sense, the logic of a modal
operator over that complement.

\section{Limitations and release recommendations}
\label{sec:limitations}

The mathematical result is unconditional within the stated semantics, but a
responsible release should keep the following boundaries visible.

\begin{itemize}
  \item The main theorem is concept satisfiability for the specified role and
  syntax, not ontology reasoning with general axioms.

  \item The direct Lean representation uses arbitrary nonempty sets.  The
  regular-closed corollary relies on the external representation theorem and
  concerns substructures, not full topological region structures.

  \item The explicit algorithm is prohibitively expensive.  It should not be
  advertised as an implemented practical reasoner, and no optimal complexity
  result is claimed.

  \item The repository is cumulative and contains superseded proof routes.
  The release should designate the cone-scheme capstones as normative, move
  obsolete theorem-level narratives out of the main path, and include the
  exact theorem-to-source map from \cref{tab:theoremmap}.

  \item A pinned Lean toolchain, a minimal build target, archived build output,
  and a stable commit or DOI should accompany the preprint.

  \item Independent human review by researchers in description logic,
  qualitative spatial reasoning, and Lean remains desirable.  Machine
  checking removes proof-script gaps relative to the definitions; it does not
  by itself establish that the definitions match every convention a reader
  may associate with the name $\Logic$.
\end{itemize}

Promising follow-up work within the decided fragment includes a symbolic
implementation, an exact complexity classification, extension to restricted
TBoxes, and a Lean formalization of the regular-closed representation bridge.

\section{Conclusion}

Removing universal partial-overlap restrictions after negation normalization
leaves a substantial logic: existential overlap, full disjointness, both
proper-part directions, arbitrary Boolean structure, and models with infinite
chains all remain.  Its satisfiability problem is nevertheless decidable.

The proof separates finiteness from model size.  A finite cone scheme records
only closure types and their possible strict predecessor types.  Any genuine
model yields a post-fixed point, which proves completeness without quotienting
the model.  Conversely, a surviving scheme is unfolded into fresh
occurrences.  Their vertical births and disjointness births form an
ordered-disjoint structure, so RCC5 closure holds globally and existential
demands never need to reuse a node.  The only missing truth-lemma branch is
$\forall\PO$, exactly the constructor excluded by the fragment.

This establishes a machine-checked positive result while sharpening the open
full-logic problem.  Ordered-disjoint normal form is the semantic core;
controlling the residual $\PO$ neighbourhoods is the remaining frontier.

\section*{Artifact availability and authorship transparency}

The formalization, probes, review packets, and historical development are
available at \url{https://github.com/lambdamikel/alcircc5}; this paper targets
the commit identified in \cref{sec:formalization}.  The $\Logic$ research
programme and split-forest idea originate with Michael Wessel.  The 2026
project used multiple AI systems under his direction for proposing proof
architectures, producing Lean code, testing, and adversarial review.  In
particular, GPT-5.6 Sol proposed the finite-control/fresh-occurrence cone-scheme
architecture, and Claude systems carried out the principal Lean development.
This manuscript was prepared with OpenAI Codex from the audited repository
state.  AI systems are acknowledged as tools rather than listed as authors;
the human author remains responsible for checking the text, claims, and public
release.

\begin{thebibliography}{99}

\bibitem{BaaderEtAl2007}
F.~Baader, D.~Calvanese, D.~L. McGuinness, D.~Nardi, and P.~F.
Patel-Schneider, editors.
\newblock \emph{The Description Logic Handbook: Theory, Implementation, and
Applications}, second edition.
\newblock Cambridge University Press, 2007.

\bibitem{BaaderRydval2020}
F.~Baader and M.~Rydval.
\newblock Description logics with concrete domains and general concept
inclusions revisited.
\newblock In \emph{Proceedings of IJCAR 2020}, LNCS 12166, pages 413--431.
Springer, 2020.

\bibitem{Cohn1993}
A.~G. Cohn.
\newblock Modal and non-modal qualitative spatial logics.
\newblock In \emph{Proceedings of the IJCAI Workshop on Spatial and Temporal
Reasoning}, 1993.

\bibitem{deMouraUllrich2021}
L.~de Moura and S.~Ullrich.
\newblock The Lean~4 theorem prover and programming language.
\newblock In A.~Platzer and G.~Sutcliffe, editors, \emph{Automated Deduction---CADE~28},
LNCS 12699, pages 625--635. Springer, 2021.
\newblock \url{https://doi.org/10.1007/978-3-030-79876-5_37}.

\bibitem{LutzMilicic2007}
C.~Lutz and M.~Mili\v{c}i\'c.
\newblock A tableau algorithm for description logics with concrete domains and
general TBoxes.
\newblock \emph{Journal of Automated Reasoning}, 38:227--259, 2007.

\bibitem{LutzWolter2006}
C.~Lutz and F.~Wolter.
\newblock Modal logics of topological relations.
\newblock \emph{Logical Methods in Computer Science}, 2(2:5):1--41, 2006.
\newblock \url{https://doi.org/10.2168/LMCS-2(2:5)2006}.

\bibitem{RandellCuiCohn1992}
D.~A. Randell, Z.~Cui, and A.~G. Cohn.
\newblock A spatial logic based on regions and connection.
\newblock In \emph{Proceedings of the Third International Conference on
Principles of Knowledge Representation and Reasoning (KR'92)}, pages 165--176,
1992.

\bibitem{RenzNebel1999}
J.~Renz and B.~Nebel.
\newblock On the complexity of qualitative spatial reasoning: A maximal
tractable fragment of the Region Connection Calculus.
\newblock \emph{Artificial Intelligence}, 108(1--2):69--123, 1999.

\bibitem{Wessel2003}
M.~Wessel.
\newblock Qualitative spatial reasoning with the $\mathcal{ALCI}_{\mathrm{RCC}}$
family---first results and unanswered questions.
\newblock Technical Report FBI-HH-M-324/03, Department of Computer Science,
University of Hamburg, May 2002 / October 2003.
\newblock \url{https://github.com/lambdamikel/alcircc5/blob/master/papers/report7.pdf}.

\bibitem{Artifact2026}
M.~Wessel.
\newblock $\Logic$ decidability project: Lean formalization and verification
artifact.
\newblock GitHub repository, commit
\code{795b2707b72e927c56b465ef0413b77aeb9a66a6}, 2026.
\newblock AI-assisted formalization and review.
\newblock \url{https://github.com/lambdamikel/alcircc5}.

\end{thebibliography}

\appendix

\section{Formal capstone signatures}
\label{app:signatures}

For reference, the central Lean interfaces in the audited snapshot are:

\begin{lstlisting}
def POFree : Concept -> Prop

theorem coneScheme_complete
  (hI : RCC5Interp I) (C0 : Concept) {x : alpha}
  (hx : I.dom x) (hC0 : sat I x C0) (n : Nat) :
  exists q in gfpIter pruneSig (sigStatic C0) n,
    C0 in q.1

theorem coneScheme_sound
  (hXS  : forall q in X, q in sigStatic C0)
  (hfix : forall q in X, q in pruneSig X)
  (hpo  : POFree C0) (hq0 : q0 in X)
  (hC0  : C0 in q0.1) :
  Satisfiable C0

theorem coneScheme_correct_at (C0 : Concept)
  (hpo : POFree C0) :
  Satisfiable C0 <->
    exists q in
      gfpIter pruneSig (sigStatic C0) (sigStatic C0).length,
      C0 in q.1

def decidableSat_cone (C0 : Concept)
  (hpo : POFree C0) : Decidable (Satisfiable C0)

def decidableFSat (F : Formula)
  (h : FPOFree true F) : Decidable (FSatisfiable F)

theorem satisfiable_iff_set (C0 : Concept) :
  Satisfiable C0 <-> SetSatisfiable C0

def decidableSetSat (C0 : Concept)
  (hpo : POFree C0) : Decidable (SetSatisfiable C0)
\end{lstlisting}

\end{document}
